\chapter{Tradycyjne Zupy Wigilijne}

\section{Barszcz czerwony z uszkami}
Wigilia w wielu regionach Polski, zwłaszcza na południu i wschodzie, tradycyjnie rozpoczyna się od barszczu czerwonego\index{Barszcz}. Zupa ta serwowana jest z drobnymi pierożkami faszerowanymi grzybami leśnymi, zwanymi uszkami\index{Uszka}. Opisuje to szczegółowo A. Nowak \cite{nowak2018}.

\section{Zupa grzybowa}
W centralnej i północnej części Polski na stołach częściej gości zupa grzybowa\index{Zupa grzybowa}, przygotowywana z suszonych borowików i podawana z domowym makaronem (tzw. łazankami).